\subsection{积分作为参数函数的连续性}

命题 6（IV，第 171 页）表明：当一个微分方程

\[
\mathbf {x} ^ {\prime} = \mathbf {f} (t, \mathbf {x})
\]

“接近”于一个利普希茨方程 \(\mathbf{x}^{\prime} = \mathbf{g}(t, \mathbf{x})\) ，且当假定两个方程在同一区间上都有近似解时，这两个近似解是“接近”的；我们将阐明这个结果，办法是证明：只要利普希茨方程 \(\mathbf{x}^{\prime} = \mathbf{g}(t, \mathbf{x})\) 在某区间上存在解，便蕴含方程 \(\mathbf{x}^{\prime} = \mathbf{f}(t, \mathbf{x})\) 在同一区间上近似解的存在，前提是在该区间上，方程 \(\mathbf{x}^{\prime} = \mathbf{g}(t, \mathbf{x})\) 的解的值并不“太接近” H 的边界。

命题 7. 设 f 与 g 是定义在 I × H 上的两个函数，满足 IV 第 165 页引理 1 的假设，并且在 I × H 上

\[
\left\| \mathbf {f} (t, \mathbf {x}) - \mathbf {g} (t, \mathbf {x}) \right\| \leqslant \alpha .\tag{16}
\]

再设 \(\mathbf{g}\) 关于常数 \(k > 0\) 在 \(\mathrm{I} \times \mathrm{H}\) 上是利普希茨的，且 \(\mathbf{f}\) 在 \(\mathrm{I} \times \mathrm{H}\) 上是局部利普希茨的，或者 \(\mathrm{E}\) 是有限维的。设 \((t_0, \mathbf{x}_0)\) 是 \(\mathrm{I} \times \mathrm{H}\) 的一点，\(\mu\) 是一个数 \(>0\) ，并设

\[
\varphi (t) = \mu e ^ {k (t - t _ {0})} + \alpha \frac {e ^ {k (t - t _ {0})} - 1}{k}.
\]

设 \(\mathbf{u}\) 是方程 \(\mathbf{x}' = \mathbf{g}(t, \mathbf{x})\) 的一个积分，它定义在包含于 I 的区间 \(\mathrm{K} = [t_0, b[\) 上、等于 \(\mathbf{x}_0\) 于点 \(t_0\) ，并且对一切 \(t \in \mathrm{K}\) ，以 \(\mathbf{u}(t)\) 为中心、以 \(\varphi(t)\) 为半径的闭球都包含于 H 中。在这些条件下，对每个满足 \(\mathbf{y} \in \mathrm{H}\) 的 \(\| \mathbf{y} - \mathbf{x}_0 \| \leqslant \mu\) ，便存在 \(\mathbf{v}\) ，它是方程 \(\mathbf{x}' = \mathbf{f}(t, \mathbf{x})\) 的一个积分，定义在 K 上、取值于 H 中且等于 \(\mathbf{y}\) 于点 \(t_0\) ；此外，在 K 上有 \(\| \mathbf{u}(t) - \mathbf{v}(t) \| \leqslant \varphi(t)\) 。

设 M 是 \(\mathbf{x}^{\prime} = \mathbf{f}(t, \mathbf{x})\) 的积分所成的族，其中每个积分都取值于 H 中、在 \(t_{0}\) 处等于 y、且定义在一个包含于 I 的半开区间 \([t_{0}, \tau[\) 上（该区间随所考虑的积分而异）。由 IV 第 171 页的定理 1（当 f 是局部利普希茨时）或 IV 第 167 页的推论（当 E 是有限维时），M 非空，且与 IV 第 166 页命题 3 中同样的推理表明，M 依序关系“v 是 w 的限制”是归纳的。设 \(v_{0}\) 是 M 的一个极大元，\([t_{0}, t_{1}[\) 是 \(v_{0}\) 的定义区间；由 IV 第 171 页的命题 6，一切都归结为证明 \(t_{1} \geqslant b\) 。若不然，则有

\[
\left\| \mathbf {u} (t) - \mathbf {v} _ {0} (t) \right\| \leqslant \varphi (t)
\]

在区间 \([t_0, t_1[\) 上此式由命题 6 得出；而在紧区间 \([t_0, t_1]\) 上，规则函数 \(\mathbf{g}(t, \mathbf{u}(t))\) 有界，故函数 \(\mathbf{g}(t, \mathbf{v}_0(t))\) 在区间 \([t_0, t_1[\) 上有界，因为在此区间上 \(\| \mathbf{g}(t, \mathbf{v}_0(t)) \| \leqslant \| \mathbf{g}(t, \mathbf{u}(t)) \| + k\varphi(t)\) 。由于 \(\mathbf{v}_0\) 是 \(\mathbf{x}' = \mathbf{g}(t, \mathbf{x})\) 的精确到 \(\alpha\) 的近似解（在 \([t_0, t_1[\) 上），故存在数 \(M > 0\) 使得在此区间上除去一个可数集的点以外有 \(\| \mathbf{v}_0'(t) \| \leqslant M\) ；于是中值定理表明 \(\| \mathbf{v}_0(s) - \mathbf{v}_0(t) \| \leqslant M |s - t|\) ，对每对点 \(s\) 与 \(t\)（都属于 \([t_0, t_1[\) ）成立，从而（由柯西准则）\(\mathbf{v}_0(t)\) 有左极限 \(\mathbf{c}\) 于点 \(t_1\) ，并且由连续性有 \(\| \mathbf{c} - \mathbf{u}(t_1) \| \leqslant \varphi(t_1)\) ，故 \(\mathbf{c} \in H\) 。于是由 IV 第 171 页定理 1 或 IV 第 167 页推论可见，存在 \(\mathbf{x}' = \mathbf{f}(t, \mathbf{x})\) 的一个积分，它定义在区间 \([t_1, t_2[\) 上且等于 \(\mathbf{c}\) 于 \(t_1\) ，这与 \(\mathbf{v}_0\) 的定义相矛盾。

定理 3. 设 F 是一个拓扑空间，f 是从 I × H × F 到 E 中的一个映射，使得对每个 \(\xi \in F\) ，函数 \((t, \mathbf{x}) \mapsto \mathbf{f}(t, \mathbf{x}, \xi)\) 在 I × H 上是利普希茨的，并且当 \(\xi\) 趋于 \(\xi_{0}\) 时，\(\mathbf{f}(t, \mathbf{x}, \xi)\) 在 I × H 上一致趋于 \(\mathbf{f}(t, \mathbf{x}, \xi_{0})\) 。设 \(\mathbf{u}_{0}(t)\) 是 \(\mathbf{x}' = \mathbf{f}(t, \mathbf{x}, \xi_{0})\) 的一个积分，它定义在包含于 I 中的区间 J = {[}t₀, b{[} 上、取值于 H 中且在 t₀ 处等于 \(x_{0}\) 。对每个包含于 J 中的紧区间 {[}t₀, t₁{]} ，存在 \(\xi_{0}\) 在 F 中的一个邻域 V，使得对每个 \(\xi \in V\) ，方程 \(\mathbf{x}' = \mathbf{f}(t, \mathbf{x}, \xi)\) 有一个（且仅有一个）积分 \(\mathbf{u}(t, \xi)\) ，它定义在 {[}t₀, t₁{]} 上、取值于 H 中且在 t₀ 处等于 \(x_{0}\) ；而且，当 \(\xi\) 趋于 \(\xi_{0}\) 时，解 \(\mathbf{u}(t, \xi)\) 在 \(\mathbf{u}_{0}(t)\)[{t₀, t₁}]{ 上一致趋于 } 。

的确，取 r \textgreater{} 0，使得对 \(t_{0} \leqslant t \leqslant t_{1}\) ，以 \(\mathbf{u}_{0}(t)\) 为中心、以 r 为半径的闭球包含于 H 中；若 \(\mathbf{f}(t, \mathbf{x}, \xi)\) 关于常数 k \textgreater{} 0 在 I × H 上是利普希茨的，则取 \(\alpha\) 足够小，使 \(\alpha \frac{e^{k(t_{1}-t_{0})} - 1}{k} < r\) ；取 V 使得对每个 \(\xi \in V\) ，在 I × H 上有 \(\|f(t, \mathbf{x}, \xi) - f(t, \mathbf{x}, \xi_{0})\| \leqslant \alpha\) ，就可以应用 IV 第 174 页的命题 7；此外，

\[
\| \mathbf {u} (t, \xi) - \mathbf {u} _ {0} (t) \| \leqslant \alpha \frac {e ^ {k \left(t _ {1} - t _ {0}\right)} - 1}{k}
\]

在 \([t_{0}, t_{1}]\) 上，这就完成了定理的证明。

注. 当 H = E 且 IV 第 174 页的条件 (16) 在 I × E 上满足时，IV 第 174 页的命题 7 可应用于 \(\mathbf{x}' = \mathbf{g}(t, \mathbf{x})\) 的每个解 u，在其有定义的任何区间上；甚至还可以假定 \(\mathbf{g}(t,\mathbf{x})\) 关于一个规则函数 \(k(t)\) 是利普希茨的，尽管该函数在此区间上未必有界。

